\documentclass[12pt]{ximera}
\typeout{Start loading xmPreamble.tex}%

% Add here extra macro's that are loaded automatically by all documents of claas 'ximera' or 'xourse' in this repo

%%
%%  Example:
%%
% \newcommand{\R}{\mathbb{R}}
\usepackage{tikz}
\usetikzlibrary{positioning, arrows.meta}
\usepackage{multicol}
\usepackage{enumitem}
\usepackage{caption}
\usepackage{amsmath}
\usepackage{amsthm}

%% NOTE: do NOT add \usepackage to individual activity .tex files.
%% When a xourse inlines an activity it gobbles \documentclass and \input, but a bare
%% \usepackage still executes -- after \begin{document} -- and errors out.
%% All shared packages and macros belong here.
%%
%% ximera.cls already loads: amsmath, amssymb, amsthm, cancel, enumitem, graphicx,
%% hyperref, listings, pgfplots, tikz, url, xcolor. No need to re-load those.

\title{Rough Estimations: Ground Rules}
\begin{document}
\begin{abstract}
What a rough estimation is (and is not), why we round as we go rather than all at
once, and how to tell whether an estimate is an underestimate or an overestimate
without ever computing the exact value.
\end{abstract}
\maketitle

\section*{Rough Estimations}

This section works differently from the others:

\begin{itemize}
  \item We want \emph{estimates}, not exact values.
  \item They are \emph{rough} --- round as much as you need to make the calculation easy.
  \item No calculators, just for this section.
  \item We use the \emph{structure} of the problem to decide which operations to perform.
  \item We round numbers so those operations can be done mentally, without long-form
    calculation.
  \item We round \emph{as we go}, not all at once.
\end{itemize}

\begin{example}
Roughly estimate how many centimeters are in $3$ miles.

We will use the equivalences
\[
12 \text{ inches} = 1 \text{ foot}, \qquad
2.54 \text{ cm} = 1 \text{ in}, \qquad
5280 \text{ ft} = 1 \text{ mile}.
\]

\textbf{Structure of the answer.} First decide what the calculation looks like, before
rounding anything:
\[
3 \cancel{\text{miles}} \cdot \frac{5280 \cancel{\text{ft}}}{1 \cancel{\text{mile}}}
\cdot \frac{12 \cancel{\text{in}}}{1 \cancel{\text{ft}}}
\cdot \frac{2.54 \text{ cm}}{1 \cancel{\text{in}}}.
\]

\textbf{Now round, one step at a time.}
\[
\begin{aligned}
3 \cancel{\text{miles}} \cdot \frac{5000 \text{ ft}}{1 \cancel{\text{mile}}}
  &= 15{,}000 \text{ ft} \\
15{,}000 \cancel{\text{ft}} \cdot \frac{10 \text{ in}}{1 \cancel{\text{ft}}}
  &= 150{,}000 \text{ in} \\
150{,}000 \cancel{\text{in}} \cdot \frac{2 \text{ cm}}{1 \cancel{\text{in}}}
  &= 300{,}000 \text{ cm}.
\end{aligned}
\]

There are roughly $300{,}000$ cm in $3$ miles.
\end{example}

\begin{problem}
Is the estimate above an underestimate, an overestimate, or neither? Can we tell
without determining the exact value?

\begin{multipleChoice}
\choice[correct]{An underestimate.}
\choice{An overestimate.}
\choice{Neither --- there is no way to tell without computing the exact value.}
\end{multipleChoice}

\begin{solution}
We rounded $5280$ down to $5000$ and multiplied by it; we rounded $12$ down to $10$
and multiplied by it; we rounded $2.54$ down to $2$ and multiplied by it. Every
rounding made a factor \emph{smaller}, and every factor was multiplied in, so the
result must be smaller than the true value. Our rough estimate is an
\textbf{underestimate}.
\end{solution}
\end{problem}

\begin{problem}
Explain in your own words why we round \emph{as we go} instead of rounding every
number at the start.

\begin{freeResponse}
\end{freeResponse}
\end{problem}

\end{document}
